\section{Vector math}

Foundational matrix calculus + matrix-statistics rules.
Convention: $\bbeta,\xb \in \R^{p}$ are column vectors;
$\partial(\,\cdot\,)/\partial\bbeta$ is a column vector
(numerator-layout for scalar-by-vector).

% ============================================================
\subsection{Matrix derivatives (scalar by vector)}
% Source: wiki/book-deltas/m03-linreg.md Step 3 of OLS derivation,
%         wiki/concepts/least-squares-and-mle.md (RSS expansion).
% These two identities are the only matrix-calculus rules the prof
% explicitly invokes; everything else is built from them.

\begin{center}
\renewcommand{\arraystretch}{1.25}
\begin{tabularx}{\linewidth}{@{}l l X@{}}
\toprule
\textbf{Expression $f(\bbeta)$} & \textbf{$\partial f/\partial \bbeta$} & \textbf{Notes} \\
\midrule
$\bm{a}\T \bbeta = \bbeta\T \bm{a}$ & $\bm{a}$
  & Scalar; transpose of a $1{\times}1$ is itself. \\
$\bbeta\T \bm{A}\, \bbeta$ & $(\bm{A} + \bm{A}\T)\bbeta$
  & General square $\bm{A}$. \\
$\bbeta\T \bm{A}\, \bbeta$ & $2\bm{A}\bbeta$
  & \textbf{Symmetric} $\bm{A}=\bm{A}\T$ (this is the case OLS uses). \\
$\bbeta\T \bbeta = \|\bbeta\|^2$ & $2\bbeta$
  & Special case, $\bm{A}=\bm{I}$. \\
$\|\Xb\bbeta - \yb\|^2$ & $2\Xb\T\Xb\bbeta - 2\Xb\T\yb$
  & OLS RSS; sets to $\bm{0}$ $\Rightarrow$ normal eqns. \\
\bottomrule
\end{tabularx}
\end{center}

\textbf{Scalar-by-vector chain rule.} For $h(\bbeta) = g(\bm{u}(\bbeta))$
with $\bm{u}:\R^p\!\to\!\R^m$ and $g:\R^m\!\to\!\R$:
$$\tfrac{\partial h}{\partial \bbeta} \;=\;
\left(\tfrac{\partial \bm{u}}{\partial \bbeta}\right)\!\T
\tfrac{\partial g}{\partial \bm{u}}.$$
% Source: wiki/concepts/backpropagation.md — "Backprop is just the
% chain rule, applied so you don't recompute things." [L24]
% The prof's framing: chain rule = multivariate calculus, no magic.

\textbf{OLS derivation in three lines}
% Source: wiki/book-deltas/m03-linreg.md §2 Steps 1–5;
%         wiki/concepts/least-squares-and-mle.md "Formulas to know cold".
\begin{align*}
\mathrm{RSS}(\bbeta) &= (\yb-\Xb\bbeta)\T(\yb-\Xb\bbeta) \\
 &= \yb\T\yb - 2\bbeta\T\Xb\T\yb + \bbeta\T\Xb\T\Xb\,\bbeta, \\
\tfrac{\partial\mathrm{RSS}}{\partial\bbeta}
 &= -2\Xb\T\yb + 2\Xb\T\Xb\,\bbeta \;\overset{!}{=}\; \bm{0}, \\
\hat\bbeta &= (\Xb\T\Xb)^{-1}\Xb\T\yb. \quad\square
\end{align*}
The two cross-terms $\yb\T\Xb\bbeta$ and $\bbeta\T\Xb\T\yb$ collapse
because each is a $1{\times}1$ scalar and equals its own transpose.

% TODO: gap — the prof does not give an explicit identity for
% d(Ax)/dx (the Jacobian). Under numerator layout it is A;
% under denominator layout it is A^T. Not used in this course.

% ============================================================
\subsection{Expectation, variance, covariance for random vectors}
% Source: wiki/concepts/random-vector-and-covariance.md
%   "Formula(s) to know cold" + L04-statlearn-3 (prof's on-board proof).
% These are the two formulas the M2 quiz (modules/2StatLearn/2StatLearn.2.md
% Q4/Q5) drills.

Let $\bm{X}\in\R^p$ be a random vector,
$\bm{\mu}=\E[\bm{X}]$,
$\bm\Sigma=\Cov(\bm{X})=\E[(\bm{X}-\bm{\mu})(\bm{X}-\bm{\mu})\T]$.
$\bm{A},\bm{B},\bm{C}$ are constant (non-random) matrices of compatible shape.

\begin{center}
\renewcommand{\arraystretch}{1.25}
\begin{tabularx}{\linewidth}{@{}l X@{}}
\toprule
\textbf{Rule} & \textbf{Statement} \\
\midrule
Sum of expectations & $\E[\bm{X}+\bm{Y}] = \E[\bm{X}] + \E[\bm{Y}]$ \\
Linear $\E$ of matrix & $\E[\bm{A}\bm{X}\bm{B}] = \bm{A}\,\E[\bm{X}]\,\bm{B}$ \\
Affine $\E$ & $\E[\bm{A}\bm{X}+\bm{b}] = \bm{A}\bm{\mu}+\bm{b}$ \\
\addlinespace
$\bm\Sigma$ definition & $\bm\Sigma = \E[\bm{X}\bm{X}\T] - \bm{\mu}\bm{\mu}\T$ \\
\addlinespace
Linear transform of cov &
  $\Cov(\bm{C}\bm{X}) = \bm{C}\,\bm\Sigma\,\bm{C}\T$
  \quad \\
Affine transform of cov &
  $\Cov(\bm{A}\bm{X}+\bm{b}) = \bm{A}\,\bm\Sigma\,\bm{A}\T$ \\
Scalar quadratic form &
  $\Var(\bm{a}\T\bm{X}) = \bm{a}\T\bm\Sigma\,\bm{a} \ge 0$
  \;(so $\bm\Sigma$ is PSD) \\
Cross-covariance &
  $\Cov(\bm{A}\bm{X},\bm{B}\bm{Y}) = \bm{A}\,\Cov(\bm{X},\bm{Y})\,\bm{B}\T$ \\
\addlinespace
Univariate sanity & $\E[aX+b]=a\E[X]+b$,\;
   $\Var(aX+b)=a^2\Var(X)$ \\
Sum of two & $\Var(X+Y)=\Var(X)+\Var(Y)+2\Cov(X,Y)$ \\
\bottomrule
\end{tabularx}
\end{center}

% ============================================================
\subsection{Sample covariance matrix from a data matrix}
% Source: wiki/concepts/principal-component-analysis.md "Notation & setup"
%         + wiki/book-deltas/m10-unsuper.md §1 (PCA pipeline).
% Prof: column-centered (and standardized for PCA) before PCA.

Data matrix $\Xb\in\R^{n\times p}$: rows = observations, columns = variables
(ISLP convention; some books transpose).
Let $\bm{1}_n\in\R^n$ be the all-ones column,
$\bar{\bm{x}}=\tfrac{1}{n}\Xb\T\bm{1}_n\in\R^p$ the column-means vector.

\begin{enumerate}
\item \textbf{Center} (subtract column means from every row):
$$\Xb_c \;=\; \Xb - \bm{1}_n\,\bar{\bm{x}}\T \;\in\; \R^{n\times p}.$$
Now each column of $\Xb_c$ has mean $0$.

\item \textbf{(For PCA only)} Standardize: divide each column of $\Xb_c$
by its sample SD $s_j=\sqrt{\tfrac{1}{n-1}\sum_i(x_{ij}-\bar x_j)^2}$,
giving the $z$-score matrix $\Xb_z$. Then every column has mean $0$,
sample variance $1$, and $\sum_k \lambda_k = p$.
% Source: wiki/book-deltas/m10-unsuper.md (standardization).

\item \textbf{Sample covariance matrix:}
$$\boxed{\;\bm{S} \;=\; \tfrac{1}{n-1}\,\Xb_c\T \Xb_c \;\in\; \R^{p\times p}.\;}$$
Symmetric, PSD. Entries: $S_{jk}=\tfrac{1}{n-1}\sum_i (x_{ij}-\bar x_j)(x_{ik}-\bar x_k)$.
Used as $\bm\Sigma$ in PCA's eigendecomposition.
\end{enumerate}

\textbf{Connection to OLS algebra:} if $\Xb$ \emph{already includes the
intercept column of ones}, then $\Xb\T\Xb$ (no centering, no $1/(n{-}1)$)
is the un-normalized cross-product used in $\hat\bbeta=(\Xb\T\Xb)^{-1}\Xb\T\yb$
and $\Cov(\hat\bbeta)=\sigma^2(\Xb\T\Xb)^{-1}$. Centering kills the off-diagonal
between intercept and slopes — \emph{numerical} reason to center.
% Source: wiki/book-deltas/m03-linreg.md §6 "(C5) Why centering helps numerics."
